\section{Unknown local directions and nuisance}\label{unknown-local-directions-and-nuisance}

The known-tilt subfamily already realizes the hardest first-order direction around its known baseline. Define the influence function
\[
\begin{gathered}
\psi(z)=\frac{U(0,z)}{\lambda s} \\
E_0\psi^2=c^2.
\end{gathered}
\]
For any bounded centered score h with \(I_h=E_0h^2>0\), differentiating the root equation gives normalized root sensitivity \(d_h=E_0[\psi h]\). Hence
\[
\frac{d_h^2}{I_h}\le c^2,
\]
with equality precisely for a nonzero multiple of the entropy gradient \(U(0,\cdot)\). The comparison fixes the seed window, not an arbitrary common parameter window with different attainable seeds.

Consider all unknown measurable k in
\[
\begin{gathered}
\mathcal K_{B,K}=\{k:E_0k=0,\ \|k\|_\infty\le B,\ |E_0\psi k|\le K\} \\
f_{n,k}=f_0(1+k/\sqrt n).
\end{gathered}
\]
The baseline, B and K are known; k is not. For large n, these densities preserve legality with the same signed density g and the same complete risk. State derivatives of U are bounded, so the whole class lies in a common shrinking smooth-field neighborhood. A repeated stopped-seed proof gives
\[
\sup_{k\in\mathcal K_{B,K}}|\pi_n(Q_{n,k})-\Phi(E_0\psi k)|\le Ce_n.
\]

\textbf{Theorem 4 (Unknown-direction minimax value; N.3).} Suppose \(B\ge K\|\psi\|_\infty/c^2\). For \(M/n\) tending to positive finite \(\rho\), the minimax prediction risk over this unknown-direction class converges to \(V_G(K,c/\sqrt\rho)\). Uniform consistency is possible exactly when \(M/n\) tends to infinity; for \(M/n\) tending to zero the value tends to \(\Phi(K)-1/2\).

The lower bound restricts to \(k_t=t\psi/c^2\), \(|t|\le K\), which the envelope condition guarantees is contained in the class. For the upper bound use the direction-free statistic
\[
\begin{gathered}
\widehat t=\frac{\sqrt n}{M}\sum_{j=1}^M\psi(\widetilde Z_j) \\
E\widehat t=E_0\psi k \\
\operatorname{Var}_{n,k}\psi=c^2+O(n^{-1/2}).
\end{gathered}
\]
Its uniform CLT transfers any bounded Lipschitz Gaussian rule's risk. Gaussian smoothing approximates an arbitrary near-optimal bounded rule with no asymptotic loss, completing the upper bound. Fixed finite step rules also transfer because the limiting Gaussian puts vanishing probability near their thresholds.

This is equality of minimax decision values, not scalar equivalence of the full infinite-dimensional experiment. Orthogonal directions may still contain additional information about nuisance. N.4 gives, for a finite score span with information matrix I and \(d=E_0[H\psi]\), the variance \(d^\top I^{-1}d/\rho=\|\Pi_{\operatorname{span}(H)}\psi\|_2^2/\rho\le c^2/\rho\). Root-preserving nuisance is orthogonal to \(\psi\); it causes no additional first-order penalty for the influence-aligned entropy tilt. These facts do not extend without a new argument to an unknown baseline or growing local envelope. Complete proofs and projection qualifications are in N.1--N.4.

\subsection{Root sensitivity and the information bound}\label{root-sensitivity-and-the-information-bound}

Differentiate E\_epsilon U(q\_epsilon,Z)=0. At epsilon=0 the state derivative is \(\lambda\) and the distribution derivative is \(E_0\) U0 h, proving (N.2). Differentiating log f\_epsilon at zero gives score h and information \(E_0\) h\^{}2. Cauchy--Schwarz gives (\(E_0\) \(\psi\) h)\^{}2\textless=\(E_0\) \(\psi\)\^{}2 \(E_0\) h\^{}2. Its equality condition proves the last assertion. \(\square\)

Thus the original \(h_0\)=U0/\(\sigma_g\) is not merely a convenient tilt: it maximizes normalized seed displacement for a given local statistical distinguishability. The term ``least favorable'' here concerns estimation of this target, after fixing the target window; it is not an ordering of all nonlinear environmental shifts.

If \(\theta\)=u/\(\sqrt n\) and the score h is known, its limiting seed is \(d_h\) u, while the calibration information is \(\rho\) \(I_h\). At a fixed seed window t=\(d_h\) u in {[}-K,K{]}, the resulting Gaussian noise standard deviation is

\[
r_h=\frac{|d_h|}{\sqrt{\rho I_h}}\le\frac c{\sqrt\rho}.
\tag{N.4}
\]

For \(d_h\)=0 the first-order seed is zero, so this reparameterization is not applicable; a second-order analysis would be a different local problem. Comparing different h at a fixed u-window would change the attainable target window and would not be the same least-favorability comparison.

\subsection{Uniform exit over the unknown-direction class}\label{uniform-exit-over-the-unknown-direction-class}

For sufficiently large n, all these fields have one root q\_\{n,k\} on the same compact exit domain, and uniformly over k,

\[
\begin{gathered}
q_{n,k}=-\frac{E_0[U_0 k]}{\lambda\sqrt n}+O(n^{-1}) \\
\lambda_{n,k}=\lambda+O(n^{-1/2}) \\
s_{n,k}=s+O(n^{-1/2}).
\tag{N.6}
\end{gathered}
\]

Their true right risk-exit probabilities obey

\[
\begin{gathered}
\boxed{\sup_{k\in\mathcal K_{B_{\rm loc},K}}
|\pi_n(Q_{n,k})-\Phi(t(k))|\le C e_n,} \\
e_n=n^{-1/2}+((1+\log n)/n)^{1/3}.
\tag{N.7}
\end{gathered}
\]

\textbf{Proof.} For every fixed derivative order j needed by the stopped-seed proof,

\[
\|v_{n,k}-v_0\|_{C^j}
\le \frac{B_{\rm loc}}{\sqrt n}
\max_{0\le l\le j}\sup_{x,z}|\partial_x^lU(x,z)|.
\]

Thus the entire class lies in one shrinking smooth-field neighborhood. No derivative of k with respect to z is needed. The baseline derivative lower bound and endpoint drift signs persist uniformly. Since v\_\{n,k\}(0)=\(E_0\) U0 k/\(\sqrt n\), Taylor expansion at zero and the common positive derivative give q=O(n\^{}\{-1/2\}) and the root expansion (N.6). Bounded derivatives and bounded U\^{}2 give the remaining expansions. The baseline root noise is nonzero, so its uniform lower bound persists.

For completeness, the probabilistic part can be repeated with exactly the uniform constants of L.2 rather than applying a fixed-law theorem. All maps have inverse contraction and \(C_2\) conjugacy bounds in the same smooth neighborhood; their root-to-boundary distances and conjugate exit magnitudes are bounded away from zero. Bounded U and its derivative give common batch second/fourth moments and the state-difference noise bound. For the stopped normalized conjugacy Y and frozen-root seed W, the same orthogonality and Taylor estimates yield

\[
\begin{gathered}
\sup_{k,j}E|Y_j|^2\le C \\
\sup_k E A_{n,k}^{-2\tau}\le C/n \\
\sup_k E|Y_\tau-W|^2\le C(1+\log n)/n.
\end{gathered}
\]

Exit is finite for every member because f\_\{n,k\}\textgreater=\(f_0\)/2 supplies a common positive mass in each of the fixed forcing intervals. Frozen seed columns have uniformly bounded third moments and variances bounded below, so Berry--Esseen and seed anti-concentration hold uniformly. Choosing the same anti-concentration cutoff as in L.2 bounds the error of identifying the final exit sign by Ce\_n.~Finally,

\[
\sqrt n H_{n,k}(0)/s_{n,k}=t(k)+O(n^{-1/2})
\]

uniformly by (N.6) and the normalized \(C^2\) bound. The global risk-safe interval is unchanged, so this is actual risk exit without a component-transfer correction. This proves (N.7). \(\square\)

The proof uses a uniform bounded envelope and a finite number of smooth U derivatives. It does not require entropy or covering-number control of the function class of directions, since no uniform empirical process over k is estimated.

\subsection{Upper and lower minimax arguments}\label{upper-and-lower-minimax-arguments}

Here \(B_{\mathrm{loc}}\) denotes the envelope B used in the main statement, and the class has the same target window K.

Define the scalar statistic from independent calibration scores

\[
\widehat t=\frac{\sqrt n}{M}\sum_{j=1}^M\psi(\widetilde Z_j).
\tag{N.8}
\]

It is available without knowing k. Under \(Q_{n,k}\),

\[
\begin{gathered}
E\widehat t=t(k) \\
\operatorname{Var}(\psi(\widetilde Z))
=c^2+n^{-1/2}E_0[\psi^2 k]-t(k)^2/n
=c^2+O(n^{-1/2})
\tag{N.9}
\end{gathered}
\]

uniformly. Boundedness of \(\psi\) gives a uniform triangular-array CLT. For \(M/n\to\rho\) in (0,infinity), the distribution of \(\widehat t\) is uniformly approximated in Kolmogorov distance by N(\(t(k)\),\(c^2\)/\(\rho\)).

Let \(R^{\mathcal K}_{n,M}\) be the minimax expected absolute prediction loss over the declared unknown-direction class, using the exact \(\pi_n\) as the target. Let

\[
V_G(K,r)=\inf_d\sup_{|t|\le K}E|d(t+rZ)-\Phi(t)|.
\]

\textbf{Lower bound.} The unknown-direction class contains the subfamily

\[
\begin{gathered}
k_t=(t/c^2)\psi=(t/c)h_0 \\
|t|\le K,
\end{gathered}
\]

by the envelope condition of Theorem 4. Its seed coordinate is exactly t, and it is the original entropy-gradient tilt with \(\theta\)=t/(c \(\sqrt n\)). Restricting to this subfamily and even telling the predictor that this restriction holds can only reduce the minimax risk. The experiment-level Gaussian reduction already proved for that subfamily gives liminf R\^{}K\textgreater=\(V_G\)(K,c/\(\sqrt\rho\)). This invokes a full-experiment lower bound, not a lower bound inferred from a moment CLT.

\textbf{Proof of the upper bound.} Any fixed bounded Lipschitz Gaussian rule d, applied to \(\widehat t\), has uniformly convergent risk by (N.9), its uniform CLT, and (N.7). This suffices to approach \(V_G\) arbitrarily closely. To see why arbitrary measurable Gaussian rules do not create a gap, start with an epsilon-optimal rule valued in {[}0,1{]}. Smooth it by d\_a(x)=E{[}d(x+aZ'){]}, with independent standard-normal Z'. The smoothed rule is bounded Lipschitz. Jensen gives

\[
E_t|d_a(X)-\Phi(t)|
\le E_t|d(X+aZ')-\Phi(t)|.
\]

The last expectation differs from that under X\textasciitilde N(t,\(r^2\)) by at most TV(N(0,\(r^2+a^2\)),N(0,\(r^2\))), uniformly in t. This tends to zero as a-\textgreater0. Thus bounded Lipschitz rules are sufficient to approximate the Gaussian minimax value. Apply their uniform risk convergence, then let a and epsilon tend to zero. The upper and lower bounds agree.

For sufficiency when \(M/n\to\infty\), use \(\Phi(\widehat t)\). Its error is bounded uniformly by

\[
\frac1{\sqrt{2\pi}}\sqrt{\frac nM}
\sup_k\sqrt{\operatorname{Var}_{n,k}\psi}+Ce_n\to0.
\]

For necessity, a bounded-ratio subsequence contains the original nondegenerate least-favorable subfamily, to which the two-point lower bound applies. \(\square\)

This theorem equates a particular decision problem's minimax value. It does \textbf{not} assert that the entire infinite-dimensional experiment is equivalent to a scalar Gaussian observation: the full sample may contain information about orthogonal nuisance directions. The lower bound comes from a subexperiment and the upper bound from a direction-free rule.

If \(M/n\to0\), the same subfamily lower bound and the constant forecast upper bound also give \(R^{\mathcal K}_{n,M}\)-\textgreater{}\(\Phi(K)\)-1/2. This extends the earlier zero-information constant to the enlarged class, without characterizing the entire nuisance experiment.

The fixed finite step rules already computed for the Gaussian problem also transfer uniformly: their finitely many discontinuities have vanishing neighborhood probability under the limiting normal densities. Consequently the existing numerical minimax brackets are relevant to the enlarged local class without rerunning the optimizer or knowing k.

\subsection{Tangent projection and target-preserving nuisance}\label{tangent-projection-and-target-preserving-nuisance}

For a finite collection of bounded centered scores H=(h1,\ldots,hd) with nonsingular information I=\(E_0\) \(HH^\top\), let d=\(E_0\){[}H \(\psi\){]}. The usual local variance associated with estimating the seed functional in this score span is

\[
\frac1\rho d^\top I^{-1}d
=\frac1\rho\|\Pi_{\operatorname{span}(H)}\psi\|_2^2
\le\frac{c^2}{\rho}.
\tag{N.12}
\]

\textbf{Proof.} The projection has coefficients \(I^{-1}d\), so its squared L2 norm is \(d^\top I^{-1}d\). Projection contracts norm. Equivalently, an empirical mean of \(d^\top\) I\^{}\{-1\}H estimates the derivative functional in this local model with that first-order variance. \(\square\)

For a closed admissible tangent space T, the corresponding sensitivity supremum is

\[
\sup_{h\in T,\,\|h\|_2>0}
\frac{\langle\psi,h\rangle^2}{\|h\|_2^2}
=\|\Pi_T\psi\|_2^2.
\tag{N.13}
\]

Equality as an operational bounded density tilt requires that the maximizing direction be admissible and bounded, or be approximable by admissible bounded scores. Equation (N.13) alone does not supply the geometry and uniform exit proof for a new model class.

If an interest score h is accompanied by nuisance scores spanning N, and those nuisance directions preserve the seed functional to first order, then N is orthogonal to \(\psi\). The efficient interest score is \(h_{\rm eff}=h-\Pi_Nh\) and

\[
\frac{\langle\psi,h\rangle^2}{\|h_{\rm eff}\|_2^2}
=\frac{\langle\psi,h_{\rm eff}\rangle^2}{\|h_{\rm eff}\|_2^2}
\le c^2.
\tag{N.14}
\]

Thus nuisance can reduce information relative to a restricted one-direction model, but cannot exceed the full tangent variance \(c^2\) in this known-baseline first-order problem. In particular, for h=\(h_0\) proportional to \(\psi\) and root-preserving nuisance, \(h_{\rm eff}\)=\(h_0\): there is no additional first-order penalty.

If a proposed nuisance direction has \(E_0\) \(\psi\) k nonzero, it changes the target and is not a target-preserving nuisance. It must be included in \(t(k)\), as in Theorem N.3, rather than projected out while keeping the target fixed. Moreover, a bounded multidimensional parameter window can permit decision improvements beyond a scalar variance comparison; (N.12) is an information calculation, not a general finite-window minimax formula.

\textbf{Concrete nuisance example.} At the symmetric Gaussian baseline, let b(z)=cos(z)-\(E_0\) cos(Z). It is bounded, centered and even, whereas \(\psi\) and \(h_0\) are odd. Thus \(E_0\) \(\psi\) b=\(E_0\) \(h_0\) b=0. The two-direction perturbation k\_\{t,v\}=(t/c)\(h_0\)+v b has seed t, with an unknown bounded v changing other aspects of the score distribution. It stays in the declared unknown-direction class whenever \textbar t\textbar{} H\_*/c+\textbar v\textbar{} \textbar\textbar b\textbar\textbar\_infinity\textless=\(B_{\mathrm{loc}}\). This supplies a legal common-risk nuisance family with no additional first-order penalty. No new nuisance simulation is claimed in the present confirmation run.

\subsection{Finite baseline information and conditional
prediction}\label{finite-baseline-information-and-conditional-prediction}

We also allow the baseline marginal to be unknown in a fixed bounded
density neighborhood, while preserving the complete common
classification risk. An independent baseline sample of size \(L\)
selects the initialization through an empirical entropy root. The
prediction target is the future risk-exit probability conditional on
this realized initialization. A uniform actual-SGD approximation and a
two-sample expansion give the Gaussian interface \[
S=t+\frac{c}{\sqrt\ell}Z_L,\qquad
D=S+\frac{c}{\sqrt\rho}Z_C,\qquad
\text{target}=\Phi(S),
\] when \(L/n\to\ell\in(0,\infty)\) and \(M/n\to\rho\in(0,\infty)\). The
shared initialization noise changes the decision problem: its minimax
value is \(V_{\mathrm{FB}}\), rather than the original
deterministic-target Gaussian value with an added variance. We prove
matching upper and lower bounds over the unknown-baseline class. For
\(K>0\) and \(L/n\to\ell\in(0,\infty]\), consistent conditional
prediction remains possible exactly when \(M/n\to\infty\); superlinear
baseline information is unnecessary. In the root-preserving case
\(K=0\), the optimal limiting rule shrinks the raw displacement estimate
by \(M/(L+M)\), and its MAE has the closed form (F.35) in
\href{unknown_direction.md\#6-finite-baseline-information-and-conditional-prediction}{the
complete integrated proofs}. These claims concern a regular scalar
experiment and do not establish a reduction of neural-network
adaptation.

The added model-specific contribution is the connection between
sample-based initialization, conditional actual-SGD exit and the joint
baseline/calibration decision problem. The conjugacy and unstable-exit
mechanism, likelihood comparison, Gaussian conditioning, convex minimax
and posterior-median facts are classical. The result should not be
presented as a new general principle of statistical efficiency.

For finite positive ratios, writing \(a_P=c_P/\sqrt{\ell}\) and
\(b_P=c_P/\sqrt{\rho}\), the asymptotic minimax value over the declared
regular class is \(\sup_P V_{\mathrm{FB}}(K;a_P,b_P)\), where \[
\begin{gathered}V_{\mathrm{FB}}(K;a,b)=\inf_d\sup_{|t|\le K}E_t|d(S+bZ_C)-\Phi(S)|,\\ S=t+aZ_L.\end{gathered}
\] In the root-preserving subcase the feasible asymptotically optimal
forecast is \[
\widehat\pi_{\mathrm{pool}}=\Phi\!\left(\frac{M}{L+M}\widehat D_n\right).
\] The pooled observations affect only prediction; the future trajectory
still starts at the baseline-sample initialization. If \(K=0\) and
\(L/n\to\infty\), the constant forecast \(1/2\) is consistent without
target calibration. General-window numerical optimization and actual-SGD
confirmation of this finite-baseline experiment remain open.
